%%%%%%%%%%%%%%%%%%%%%%%%%%%%%%%%%%%%%%%%%%%%%%%%%%%%%%%%%%%%%%%%%%%%%%%%%%%
% Chapter: Op-Amp Basics
% Falstad examples: examples/opampbasics (generated by falstad/circuits/opampbasics.py)
%%%%%%%%%%%%%%%%%%%%%%%%%%%%%%%%%%%%%%%%%%%%%%%%%%%%%%%%%%%%%%%%%%%%%%%%%%%

% opfig=<s>: scales coordinates and component size by s, but not the text
\tikzset{opfig/.style={scale=#1, /tikz/circuitikz/bipoles/length=#1cm}}

%%%%%%%%%%%%%%%%%%%%%%%%%%%%%%%%%%%%%%%%%%%%%%%%%%%%%%%%%%%%%%%%%%%%%%%%%%%
\section{The Op-Amp in a Nutshell}
%%%%%%%%%%%%%%%%%%%%%%%%%%%%%%%%%%%%%%%%%%%%%%%%%%%%%%%%%%%%%%%%%%%%%%%%%%%
\nsl{The operational amplifier (op-amp) appears in almost every chapter of
this lecture: as a buffer behind a voltage reference, as the control
element of a voltage or current source, as the amplifier behind a sensor
bridge, in active filters and in the sample-and-hold circuit. This short
chapter gives the minimum that the following chapters need: what an op-amp
does, why it is always used with negative feedback, and the two golden
rules that make the analysis of op-amp circuits easy. With these rules we
look at the most useful op-amp circuit of all, the unity-gain buffer
(voltage follower), and at the non-inverting amplifier.
\newline
Chapter~\ref{ch:OpAmps} comes back to op-amps in detail: inverting,
summing, difference and instrumentation amplifiers, and the limits of real
op-amps. Good introductions are \cite{horowitz2015} (chapter~4) and
\cite{mancini2002}.
\newpage}

\begin{description}
\item[{\rot\bf Symbol and model}]~\\[-\lblskip]
\begin{center}
\begin{circuitikz}[opfig=1.3]
  \draw (0,0) node[op amp](OA){};
  \draw (OA.-) -- ++(-0.6,0) node[left]{$V_-$ (inverting input)};
  \draw (OA.+) -- ++(-0.6,0) node[left]{$V_+$ (non-inverting input)};
  \draw (OA.out) -- ++(0.6,0) node[right]{$\Vout = A\,(V_+ - V_-)$};
  \draw (OA.up) -- ++(0,0.3) node[vcc]{$+V_S$};
  \draw (OA.down) -- ++(0,-0.3) node[vee]{$-V_S$ or ground};
\end{circuitikz}
\end{center}
\bi
	\ite an amplifier for the \textbf{difference} of two voltages, with a huge gain $A = 10^5 \ldots 10^7$
	\ite inputs draw (almost) no current; the output can deliver some \SI{10}{\milli\ampere}
	\ite the output stays inside the supply range; supply pins are often not drawn
\ei
\nsl{The open-loop gain $A$ is so large that a difference of a few
hundred microvolts between the inputs is enough to drive the output to one
of the supply rails. Without feedback the op-amp only tells which input is
higher: it is a comparator (used in the flash ADC, chapter~\ref{ch:ADC}).
To get a well-defined linear circuit, a part of the output voltage is fed
back to the inverting input.}
\os{\newpage}

\item[{\rot\bf Negative feedback}]~\\[-\lblskip]
\bi
	\ite a part of $\Vout$ is fed back to the \textbf{inverting} input
	\ite $\Vout$ rises $\Rightarrow$ $V_-$ rises $\Rightarrow$ $V_+ - V_-$ shrinks $\Rightarrow$ $\Vout$ is pulled back
	\ite the circuit settles where $V_+ - V_- = \Vout/A \approx 0$ (\SI{50}{\micro\volt} for $\Vout=\SI{5}{\volt}$, $A=10^5$)
\ei
\satz{golden-basics}{\textbf{Golden rules} for an op-amp with negative feedback:
\begin{enumerate}
\item The output does whatever is necessary to make the voltage difference
      between the inputs zero: $V_+ = V_-$.
\item The inputs draw no current: $I_+ = I_- = 0$.
\end{enumerate}}
\nsl{Rule~1 holds only while the feedback loop works: the output must not
be stuck at a supply rail, and the feedback must go to the inverting
input. Rule~2 is a property of the op-amp itself; real inputs draw
picoamperes (CMOS) to some \SI{10}{\nano\ampere} (bipolar).}
\os{\newpage}

\item[{\rot\bf Real op-amps on a single supply}]~\\[-\lblskip]
\bi
	\ite most circuits in the lab run from a single supply (\SI{5}{\volt} or \SI{12}{\volt}) and ground
	\ite \textbf{input common-mode range}: the input voltages must stay inside it (often the supply minus \SI{1.5}{\volt})
	\ite \textbf{output swing}: a rail-to-rail output reaches the rails within some \SI{10}{\milli\volt}, older types (LM358) only within about \SI{1.5}{\volt} of the positive rail
	\ite \textbf{output current}: typically \SIrange{10}{40}{\milli\ampere}, short-circuit protected
	\ite examples: MCP6002, TLV9002 (CMOS, rail-to-rail, \SI{5}{\volt}), LM358 (bipolar, up to \SI{32}{\volt})
\ei
\nsl{The data sheet tells which voltages are allowed at the inputs and
which output voltages are possible at a given load current. These two
ranges are the first thing to check in every op-amp circuit, before any
formula: the golden rules only hold inside them.}
\os{\newpage}
\end{description}

%%%%%%%%%%%%%%%%%%%%%%%%%%%%%%%%%%%%%%%%%%%%%%%%%%%%%%%%%%%%%%%%%%%%%%%%%%%
\section{The Unity-Gain Buffer}
%%%%%%%%%%%%%%%%%%%%%%%%%%%%%%%%%%%%%%%%%%%%%%%%%%%%%%%%%%%%%%%%%%%%%%%%%%%
\nsl{A circuit is built from parts: a sensor divider, a filter, a
reference, an amplifier, a converter input. Each part is designed for
itself, with the formulas of chapter~\ref{ch:Basics}. When two parts are
connected, the second part loads the first one, and the formulas are no
longer valid: the voltage divider delivers less voltage, the filter has
another cut-off frequency. The unity-gain buffer (voltage follower)
separates the two parts. It copies the voltage from its input to its
output, takes no current from the first part and delivers the current for
the second part from the supply.\newpage}

\begin{description}
\item[{\rot\bf Every connection is a voltage divider}]~\\[-\lblskip]
\begin{center}
\begin{circuitikz}[scale=0.9, transform shape]
  \draw[dashed, rwucyandark] (-1.1,-0.4) rectangle (3.0,3.9);
  \node[font=\small, rwucyandark] at (1.0,4.25) {part A (Thevenin equivalent)};
  \draw (0,0) to[V, invert, l=$V_A$] (0,3) to[R, l=$R_\mathrm{out}$] (2.6,3) -- (3.7,3) node[circ]{} node[above]{$V_B$};
  \draw[dashed, rwucyandark] (4.4,-0.4) rectangle (6.6,3.9);
  \node[font=\small, rwucyandark] at (5.5,4.25) {part B};
  \draw (3.7,3) -- (5.2,3) to[R, l=$R_\mathrm{in}$] (5.2,0) -- (0,0);
\end{circuitikz}
\end{center}
\keyeq{Loading}{$\displaystyle V_B = V_A\,\frac{R_\mathrm{in}}{R_\mathrm{in}+R_\mathrm{out}}
\qquad\text{error} \approx \frac{R_\mathrm{out}}{R_\mathrm{in}}\ \text{ for } R_\mathrm{in}\gg R_\mathrm{out}$}
\bi
	\ite rule of thumb: $R_\mathrm{in} \geq 100\,R_\mathrm{out}$ for an error below \SI{1}{\percent}
\ei
\nsl{Part A is replaced by its Thevenin equivalent (chapter~\ref{ch:Basics}):
a voltage divider of $R_1$ and $R_2$ has $R_\mathrm{out} = R_1\parallelto R_2$,
an RC low-pass has the output impedance $R\parallelto (1/j\omega C)$. Part B is
represented by its input resistance. If the condition
$R_\mathrm{in}\gg R_\mathrm{out}$ cannot be met, a buffer is needed.}
\os{\newpage}

\item[{\rot\bf The voltage follower}]~\\[-\lblskip]
\begin{center}
\begin{circuitikz}[opfig=1.3]
  \draw (0,0) node[op amp, noinv input up](OA){};
  \draw (OA.+) to[short,-o] ++(-1,0) node[left]{$\Vin$ (from part A)};
  \draw (OA.-) -- ++(0,-0.9) -| ($(OA.out)+(0.4,0)$);
  \draw (OA.out) to[short,-o] ++(1.2,0) node[right]{$\Vout = \Vin$ (to part B)};
\end{circuitikz}
\end{center}
\bi
	\ite rule 1: $V_- = V_+$, and $V_- = \Vout$ $\Rightarrow$ $\Vout = \Vin$ (gain 1)
	\ite rule 2: input current $\approx 0$, input resistance \si{\giga\ohm} \ldots \si{\tera\ohm}: part A is \textbf{not loaded}
	\ite output resistance with feedback $r_o/(1+A)$: \SI{50}{\ohm}$/10^5 = \SI{0.5}{\milli\ohm}$: part B sees an \textbf{ideal source}
\ei
\keyeq{Unity-gain buffer}{$\Vout = \Vin \qquad R_\mathrm{in}\to\infty \qquad R_\mathrm{out}\to 0$}
\os{\newpage}

\item[{\rot\bf Example: potentiometer as a set-point}]~\\[-\lblskip]
\bi
	\ite potentiometer \SI{10}{\kilo\ohm} at \SI{5}{\volt}, wiper in the middle: $V_A = \SI{2.5}{\volt}$, $R_\mathrm{out} = \SI{5}{\kilo\ohm}\parallelto\SI{5}{\kilo\ohm} = \SI{2.5}{\kilo\ohm}$
	\ite load \SI{1}{\kilo\ohm} directly at the wiper: $V_B = \SI{2.5}{\volt}\cdot\frac{1}{1+2.5} = \SI{0.714}{\volt}$
	\ite load behind a follower: \SI{2.50}{\volt}; the op-amp delivers \SI{2.5}{\milli\ampere}
	\ite the wiper itself still has \SI{2.50}{\volt}: the follower does not load it
\ei
\falstad{buffer_pot}\par
\nsl{The op-amp in the Falstad example runs from \SI{5}{\volt} and ground,
like a CMOS rail-to-rail op-amp such as the MCP6002. With the slider of
the directly connected load you can see how large $R_L$ must be until the
wiper voltage is close to \SI{2.5}{\volt}. Exercise~\ref{ex:oab-pot} works
through this example.}
\os{\newpage}

\item[{\rot\bf Example: two RC low-passes in cascade}]~\\[-\lblskip]
\begin{center}
\begin{tikzpicture}
\begin{semilogxaxis}[width=12cm, height=5cm, xmin=10, xmax=10000, ymin=-50, ymax=3,
  xlabel={$f$ in \si{\hertz}}, ylabel={$|H|$ in \si{\deci\bel}}, grid=major,
  legend pos=south west, legend style={font=\scriptsize},
  ticklabel style={font=\scriptsize}, label style={font=\small}]
  \addplot[rwuviolet, very thick, domain=10:10000, samples=120]
    {-20*log10(1+(x/159.2)^2)};
  \addplot[rwucyan, very thick, domain=10:10000, samples=120]
    {-10*log10((1-(x/159.2)^2)^2+9*(x/159.2)^2)};
  \legend{with buffer, directly connected}
\end{semilogxaxis}
\end{tikzpicture}
\end{center}
\bi
	\ite each stage $R=\SI{10}{\kilo\ohm}$, $C=\SI{100}{\nano\farad}$: $f_c = \SI{159}{\hertz}$
	\ite with buffer: $H = H_1\cdot H_2 = \frac{1}{(1+j\omega RC)^2}$, at $f_c$: $|H| = 1/2$ (\SI{-6}{\deci\bel})
	\ite directly: $H = \frac{1}{1+3j\omega RC-(\omega RC)^2}$, at $f_c$: $|H| = 1/3$ (\SI{-9.5}{\deci\bel})
\ei
\falstad{buffer_rc_cascade}\par
\nsl{Without the buffer the second stage draws current from the capacitor
of the first stage. The transfer function is no longer the product of the
two single transfer functions, the cut-off frequency moves down and the
corner becomes softer. The Falstad example drives both versions with a
sine of \SI{1}{\volt} amplitude at $f_c$: the scope shows
\SI{0.333}{\volt} for the direct and \SI{0.500}{\volt} for the buffered
cascade. With the buffer each stage can be designed on its own, which is
the idea behind active filters (chapter~\ref{ch:Filters}).}
\os{\newpage}

\item[{\rot\bf Where buffers are used -- and their limits}]~\\[-\lblskip]
\bi
	\ite in this lecture: behind a voltage reference (chapter~\ref{ch:References}), at bridge taps (chapter~\ref{ch:Bridges}), between filter stages (chapter~\ref{ch:Filters}), behind the hold capacitor (chapter~\ref{ch:SampleHold})
	\ite limits
	\itee output current (some \SI{10}{\milli\ampere}) and output swing
	\itee input common-mode range: e.g.\ LM358 inputs only up to $V_S - \SI{1.5}{\volt}$
	\itee capacitive loads (long cables): the follower may oscillate; add \SIrange{50}{100}{\ohm} in series with the output
	\itee offset voltage: $\Vout = \Vin + V_\mathrm{OS}$ (\si{\micro\volt} \ldots \si{\milli\volt})
\ei
\os{\newpage}
\end{description}

%%%%%%%%%%%%%%%%%%%%%%%%%%%%%%%%%%%%%%%%%%%%%%%%%%%%%%%%%%%%%%%%%%%%%%%%%%%
\section{The Non-Inverting Amplifier}
%%%%%%%%%%%%%%%%%%%%%%%%%%%%%%%%%%%%%%%%%%%%%%%%%%%%%%%%%%%%%%%%%%%%%%%%%%%
\nsl{If only a part of the output voltage is fed back, the follower
becomes an amplifier. It keeps the good input property of the follower:
the source is not loaded. The next chapter uses it to scale a reference
voltage into an exact output voltage.\newpage}

\begin{description}
\item[{\rot\bf Non-inverting amplifier}]~\\[-\lblskip]
\begin{center}
\begin{circuitikz}[opfig=1.3]
  \draw (0,0) node[op amp, noinv input up](OA){};
  \coordinate (X) at ($(OA.out)+(0.5,0)$);
  \draw (OA.+) to[short,-o] ++(-1,0) node[left]{$\Vin$};
  \draw (OA.-) -- ++(0,-1) coordinate(N);
  \draw (N) to[R, l_=$R_2$] (N -| X) -- (X);
  \draw (N) to[R, l_=$R_1$] ++(0,-1.6) node[ground]{};
  \draw (OA.out) -- (X) to[short,-o] ++(0.8,0) node[right]{$\Vout$};
  \draw (N) node[circ]{} node[left]{$V_-$};
\end{circuitikz}
\end{center}
\bi
	\ite $R_1, R_2$ divide the output: $V_- = \Vout\,\frac{R_1}{R_1+R_2}$; rule 1: $V_- = \Vin$
	\ite $R_2 = 0$ (or $R_1\to\infty$): the follower with $G=1$
\ei
\keyeq{Non-inverting amplifier}{$\displaystyle G = \frac{\Vout}{\Vin} = 1 + \frac{R_2}{R_1}$
\quad input resistance $\approx\infty$, output resistance $\approx 0$}
\nsl{The gain depends only on the ratio of two resistors, not on the
op-amp. With \SI{1}{\percent} resistors the gain is accurate to about
\SI{2}{\percent}; with \SI{0.1}{\percent} resistors to \SI{0.2}{\percent}.
Chapter~\ref{ch:OpAmps} derives the inverting amplifier and shows how a
finite open-loop gain reduces the accuracy.}
\os{\newpage}
\end{description}

\exercise{Buffering a potentiometer}\label{ex:oab-pot}
A \SI{10}{\kilo\ohm} potentiometer between \SI{5}{\volt} and ground sets a voltage.
The wiper is in the middle position. The voltage is used by a following stage
with an input resistance of $R_L = \SI{1}{\kilo\ohm}$.
\begin{talist}
\ta Draw the Thevenin equivalent of the wiper. In which wiper position is $R_\mathrm{out}$ largest?
\ta Calculate the voltage at $R_L$ without a buffer and the error in percent.
\ta How large must $R_L$ be for an error below \SI{1}{\percent} without a buffer?
\ta Now a voltage follower (op-amp at \SI{5}{\volt}, rail-to-rail) is placed between wiper and $R_L$.
    Calculate the voltage at $R_L$, the current from the wiper and the output current of the op-amp.
\ta The wiper is turned to the upper end (\SI{5}{\volt}). What does an LM358 at \SI{5}{\volt} supply deliver, and why?
\ta Check (b) and (d) with Falstad.
\end{talist}

\begin{solution}
\begin{talist}
\ta $V_A = \SI{2.5}{\volt}$, $R_\mathrm{out} = \SI{5}{\kilo\ohm}\parallelto\SI{5}{\kilo\ohm} = \SI{2.5}{\kilo\ohm}$.
    For a wiper position $x$ (0 \ldots 1): $R_\mathrm{out} = x(1-x)\,\SI{10}{\kilo\ohm}$, largest in the middle ($R/4$).
\ta $V_B = \SI{2.5}{\volt}\cdot\SI{1}{\kilo\ohm}/(\SI{1}{\kilo\ohm}+\SI{2.5}{\kilo\ohm}) = \SI{0.714}{\volt}$, error $-\SI{71}{\percent}$.
\ta $R_L \geq 100\cdot\SI{2.5}{\kilo\ohm} = \SI{250}{\kilo\ohm}$.
\ta $V_B = \SI{2.50}{\volt}$; wiper current $\approx 0$ (picoamperes for a CMOS op-amp);
    op-amp output current $\SI{2.5}{\volt}/\SI{1}{\kilo\ohm} = \SI{2.5}{\milli\ampere}$, which comes from the \SI{5}{\volt} supply.
\ta The input common-mode range of the LM358 ends at about $V_S-\SI{1.5}{\volt} = \SI{3.5}{\volt}$, and its output
    does not reach the positive rail either: the output stays at about \SI{3.5}{\volt}. For signals up to the supply
    a rail-to-rail op-amp is needed.
\ta \falstad{buffer_pot}\newline \texttt{V\_direct} $=\SI{0.714}{\volt}$, \texttt{V\_wiper} $=$ \texttt{V\_buffered} $=\SI{2.50}{\volt}$.
\end{talist}
\end{solution}

%%% Local Variables:
%%% mode: latex
%%% TeX-master: "docu"
%%% End:
